% 图片文件等放在source目录下，需要swjtu_logo_fig_new.eps

\documentclass[12pt]{ctexbeamer}	%声明文档类型为beamer幻灯片
\mode <presentation>
\usetheme{Madrid}%%Warsaw [secheader]
%\usecolortheme{beaver}
\beamersetaveragebackground{black!10}

\setbeamercovered{transparent}
%\setbeamertemplate{navigation symbols}{}

\usefonttheme{professionalfonts}
\useinnertheme{circles}%{rectangles}
\setbeamertemplate{itemize item}{$\circledast$}%{\checkmark}


\setbeamertemplate{title page}
{
    % \vbox{}
    \vfill
    \begin{centering}
    \includegraphics[height=2.0cm]{source/swjtu_logo_fig_new.eps}
        \vskip0.6em\par%
    \begin{beamercolorbox}[sep=8pt,center,shadow=true,rounded=true]{title}
        \usebeamerfont{title}\inserttitle\par%
        \ifx\insertsubtitle\@empty%
        \else%
            \vskip0.25em%
            {\usebeamerfont{subtitle}\usebeamercolor[fg]{subtitle}\insertsubtitle\par}%
        \fi%
    \end{beamercolorbox}%
    \vskip1em\par
    \begin{beamercolorbox}[sep=8pt,center]{author}
        \usebeamerfont{author}
        % \insertauthor
        {
            \begin{tabular}{cc}
    答辩人：&\insertauthor\\
    导\quad 师：& \advisor
    \end{tabular}
        }
    \end{beamercolorbox}
    \begin{beamercolorbox}[sep=8pt,center]{institute}
        \usebeamerfont{institute}\insertinstitute
    \end{beamercolorbox}
    \begin{beamercolorbox}[sep=8pt,center]{date}
        \usebeamerfont{date}\insertdate
    \end{beamercolorbox}\vskip0.5em
    % {\usebeamercolor[fg]{titlegraphic}\inserttitlegraphic\par}
    \end{centering}
    \vfill
}

\hypersetup{pdfpagemode=FullScreen} % makes your presentation go automatically to full screen

\usepackage{fontspec,xunicode,xltxtra}

\usepackage{multimedia}	%让文档支持多媒体
\usepackage{graphics}	%让文档支持图片
\usepackage{hyperref}	%让文档支持超链接
\usepackage{booktabs}	%让文档支持三线表格
\usepackage{amsmath}	%ams可以让文档支持数学公式
\usepackage{amsfonts}
\usepackage{amssymb}
\usepackage{color}
\usepackage{graphicx,psfrag}
\usepackage{epsfig}

\title{固体潮汐形变研究}	%幻灯片标题
\author[陈昌福]{陈昌福}
\def\advisor{教授}

\institute[西南交通大学]{
大地测量学实验室\\
西南交通大学，测绘遥感信息系
}   %作者单位
\date{\today}

\setcounter{tocdepth}{1}

\begin{document}
\begin{frame}
    \titlepage
\end{frame}

\AtBeginSection[]{
    \frame{
        \frametitle{目录}
        \tableofcontents[currentsection]
    }
}

\begin{frame}
    \frametitle{目录}
    \tableofcontents
\end{frame}

%%%%%%%%%%%%%%%%%%%%%%%%%%%%%%%%%%%%%%%%%%%%%%%%%%%%%%%%%%%%%%%%%%%%%%%%%%%%%%%

\section{概述}

\begin{frame}
    \frametitle{概述}
    \begin{itemize}
        \item<1-> 研究背景
        \begin{itemize}
            \item<1-> 背景
            \item<1-> 背景
        \end{itemize}
        \item<2-> 研究意义
        \begin{itemize}
            \item<2-> 意义
            \item<2-> 意义
        \end{itemize}
    \end{itemize}
\end{frame}

%%%%%%%%%%%%%%%%%%%%%%%%%%%%%%%%%%%%%%%%%%%%%%%%%%%%%%%%%%%%%%%%%%%%%%%%%%%%%%%

\section{理论基础}

\begin{frame}
    \frametitle{基本概念}
    \begin{itemize}
        \item<1-> 定义
        \item<2-> \alert{定理}
    \end{itemize}
\end{frame}

%%%%%%%%%%%%%%%%%%%%%%%%%%%%%%%%%%%%%%%%%%%%%%%%%%%%%%%%%%%%%%%%%%%%%%%%%%%%%%%

\begin{frame}
    \frametitle{重要理论}
    \begin{itemize}
        \item<1-> 理论一
        \item<2-> 理论二
    \end{itemize}
\end{frame}

%%%%%%%%%%%%%%%%%%%%%%%%%%%%%%%%%%%%%%%%%%%%%%%%%%%%%%%%%%%%%%%%%%%%%%%%%%%%%%%

\section{研究方法}

\begin{frame}
    \frametitle{方法一}
    
    \begin{itemize}
        \item<1-> 方法边界
        \item<2-> 实现步骤
    \end{itemize}
\end{frame}

%%%%%%%%%%%%%%%%%%%%%%%%%%%%%%%%%%%%%%%%%%%%%%%%%%%%%%%%%%%%%%%%%%%%%%%%%%%%%%%

\begin{frame}
\frametitle{方法二}
$$M=m_{ubm}+Dy$$
\end{frame}

%%%%%%%%%%%%%%%%%%%%%%%%%%%%%%%%%%%%%%%%%%%%%%%%%%%%%%%%%%%%%%%%%%%%%%%%%%%%%%%

\begin{frame}
    \frametitle{方法对比}
    \begin{block}{优势}
        \begin{itemize}
            \item \alert{方法一} 基础理论
            \item \alert{方法二} 基本实现
        \end{itemize}
    \end{block}
    \pause
    \begin{block}{劣势}
        \begin{itemize}
            \item \alert{方法一} 基础理论
            \item \alert{方法二} 基本实现
        \end{itemize}
    \end{block}
\end{frame}

%%%%%%%%%%%%%%%%%%%%%%%%%%%%%%%%%%%%%%%%%%%%%%%%%%%%%%%%%%%%%%%%%%%%%%%%%%%%%%%

\section{总结与展望}

\begin{frame}
    \frametitle{实验结果}
    \begin{columns}
        \begin{column}{0.50\textwidth}
            \epsfig{figure=source/swjtu_logo_fig_new,width=\textwidth}
        \end{column}
        \begin{column}{0.50\textwidth}
            \begin{block}{研究结论}
                \begin{itemize}
                    \item 实验结果如左图所示
                    \item 实验结果如左图所示
                    \item 实验结果如左图所示
                    \item 实验结果如左图所示
                \end{itemize}
            \end{block}
        \end{column}
    \end{columns}
\end{frame}


%%%%%%%%%%%%%%%%%%%%%%%%%%%%%%%%%%%%%%%%%%%%%%%%%%%%%%%%%%%%%%%%%%%%%%%%%%%%%%%


\begin{frame}
    \centerline{\Large 欢迎大家批评指正，谢谢！}
\end{frame}

\end{document}
